% from ega2-1-fr.tex
\oldpage[II]{140}

points $a$, $b$~: $a$ est l'unique point \emph{fermé}, on a
$\mathfrak{j}_a=\mathfrak{m}$ et $k(a)=k$ est le \emph{corps résiduel}
$k=A/\mathfrak{m}$~; $b$ est le \emph{point générique} de $Y$,
$\mathfrak{j}_b=(0)$, l'ensemble $\{b\}$ étant l'unique ensemble ouvert de
$Y$ distinct de $\varnothing$ et de $Y$ (ensemble ouvert qui est donc
\emph{partout dense}), et $k(b)=K$ est le \emph{corps des fractions} de $A$.
\end{env}

\begin{env}[7.1.6]
\label{II.7.1.6-fr}
Pour un anneau local $A$, noethérien et de dimension $1$, on sait que les
conditions suivantes sont équivalentes
(\cite[p.~2-08 et 17-01]{I-1})~:
\begin{enumerate}
\item[\rm{a)}] $A$ est normal~;
\item[\rm{b)}] $A$ est régulier~;
\item[\rm{c)}] $A$ est un anneau de valuation~;
\end{enumerate}
en outre, $A$ est alors un \emph{anneau de valuation discrète}. Les prop.
(\hyperref[II.7.1.2-fr]{7.1.2}) et
(\hyperref[II.7.1.3-fr]{7.1.3}) ont alors les analogues suivants pour les
anneaux de valuation discrète~:
\end{env}

\begin{proposition}[7.1.7]
\label{II.7.1.7-fr}
Soient $A$ un anneau local intègre noethérien qui n'est pas un corps, $K$ son
corps des fractions, $L$ une extension de type fini de $K$~; il existe alors
un anneau de valuation discrète ayant $L$ pour corps des fractions et
dominant $A$.
\end{proposition}

Supposons d'abord $L=K$. Soient $\mathfrak{m}$ l'idéal maximal de $A$,
$(x_1,\ldots,x_n)$ un système de générateurs non nuls de $\mathfrak{m}$,
$B$ le sous-anneau
$A[x_2/x_1,\ldots,x_n/x_1]$ de $K$, qui est noethérien. Il est immédiat que
l'idéal $\mathfrak{m}B$ de $B$ est identique à l'idéal principal $x_1B$~;
si $\mathfrak{p}$ est un idéal premier minimal de $x_1B$, on sait que
$\mathfrak{p}$ est de rang $1$ (\cite[t.~I, p.~277]{I-13})~; autrement dit
$B_{\mathfrak{p}}$ est un anneau local noethérien \emph{de dimension $1$}~;
il est clair que
$\mathfrak{p}B_{\mathfrak{p}}\cap A$ est un idéal de $A$ contenant
$\mathfrak{m}$ et ne contenant pas $1$, donc égal à $\mathfrak{m}$, et par
suite $B_{\mathfrak{p}}$ \emph{domine} $A$
(\hyperref[I.8.1.1-fr]{I, 8.1.1}). Il résulte du th. de Krull-Akizuki
(\cite[p.~293]{II-25}) que la clôture intégrale $C$ de $B_{\mathfrak{p}}$
est un anneau noethérien (bien que $C$ ne soit pas nécessairement un
$B_{\mathfrak{p}}$-module de type fini)~; si $\mathfrak{n}$ est un idéal
maximal de $C$, $C_{\mathfrak{n}}$ est un anneau local noethérien normal de
dimension $1$ (\cite[p.~295]{II-25}), donc un anneau de valuation discrète,
qui domine $B_{\mathfrak{p}}$ et \emph{a fortiori} $A$.

Si maintenant $L$ est une extension de type fini de $K$, on peut, d'après ce
qui précède, se limiter au cas où $A$ est déjà un anneau de valuation
discrète. Soit $w$ une valuation de $K$ associée à $A$~; il existe une
valuation discrète $w'$ de $L$ qui \emph{prolonge} $w$~: on se ramène en
effet par récurrence sur le nombre de générateurs de $L$ au cas où
$L=K(\alpha)$, et alors la proposition est classique
(\cite[p.~106]{II-24}).

\begin{corollary}[7.1.8]
\label{II.7.1.8-fr}
Soient $A$ un anneau intègre noethérien, $K$ son corps des fractions, $L$ une
extension de type fini de $K$. Alors la fermeture intégrale de $A$ dans $L$
est l'intersection des anneaux de valuation discrète ayant $L$ pour corps des
fractions et contenant $A$.
\end{corollary}

En effet, un tel anneau de valuation discrète, étant normal, contient
\emph{a fortiori} tout élément de $L$ entier sur $A$. Il suffit donc de
prouver que si $x\in L$ n'est pas entier sur $A$, il existe un anneau de
valuation discrète $C$ ayant $L$ pour corps des fractions, contenant $A$ et
ne contenant pas $x$. L'hypothèse sur $x$ signifie en effet que l'on a
$x\notin B=A[1/x]$, autrement dit $1/x$ n'est pas inversible dans l'anneau
noethérien $B$. Il y a donc un idéal premier $\mathfrak{p}$ de $B$ contenant
$1/x$. L'anneau local intègre $B_{\mathfrak{p}}$ est noethérien et contenu
dans $L$, qui est une extension de type fini du corps des fractions de
$B_{\mathfrak{p}}$ (ce dernier contenant $K$). En vertu de
(\hyperref[II.7.1.7-fr]{7.1.7}), il y a donc un anneau de valuation discrète
$C$, dominant $B_{\mathfrak{p}}$ et ayant $L$ pour corps des fractions~;
comme $1/x\in\mathfrak{p}B_{\mathfrak{p}}$ appartient à l'idéal maximal de
$C$, on a $x\notin C$, ce qui conclut la démonstration.

La forme géométrique de (\hyperref[II.7.1.7-fr]{7.1.7}) est la suivante~:

